% Transcription — fonds Alexandre Grothendieck, shelfmark 43, batch 8 (pages 141-160).
% Established from the facsimile archives/batches/43/batch-08.pdf (PDF page k =
% archivists' page 140+k, no cover sheet), cut from a copy of 43.pdf fetched
% through the project's own relay
% (https://grothendieck-relay.onrender.com/source/43.pdf); tiles in
% archives/tiles/43/p141-p160. Per the skill transcribe-grothendieck. The folder
% is 207 pages in eleven batches, transcribed in parallel; this is the last
% one written. Batch 6 (ending on the EGA III leaf of p. 120), batch 7 (curves,
% ending p. 139 on Ext_{U/k}(G, gamma~) ≃ Ext_k(N(G), gamma) « ??? ») and
% batch 9 (starting p. 161 on « Tableau des Hom », with the cover of p. 170
% « Dualité de Cartier des groupes affines et groupes formels ») were read and
% their conventions followed: \underline{\mathrm{Ext}} for his underlined Ext,
% cover sheets as unnumbered section headings with a \note, published texts
% in another hand given one French \note per leaf.
%
% Pass: Opus 5.5 (claude-opus-5-5), 2026-09-26 — first pass, unchecked against the pages by a human.
%
% Dating: the folder has its own inventory date, carried verbatim in \dating{}
% with no group sentence (the group « Jacobiennes » (43-44) is added by the
% skill only for undated folders). No page of this batch carries a date; the
% Tate manuscript of pp. 142-144 cites « S. Shatz (Harvard Thesis, June
% 1962) », said in the \note of p. 143.
%
% Copies. Page 160 (grey ground, uniformly grey pencil-like strokes, like
% pp. 161-169 of batch 9) looks like a photocopy of manuscript notes in his
% hand, and is transcribed as an original. Pages 147-155 and 157 are blue ink
% on original paper (rust marks of a staple at the top left); 145 is pencil
% on tan paper; the covers 141, 146, 159 are originals. Pages 142-144 are a
% fair copy in blue-violet ink in a hand that is not his; whether they are an
% original or a duplicator copy was not established.
%
% Pages skipped: 156 and 158, blank sheets (158 carries only a stray pencil
% mark).
%
% Covers: 141 « Dualité de Tate-Cassels », 146 « Groupes de Galois comme
% groupes analytiques sur Q_l, d'après Serre », 159 « Tableau des Hom, Hom
% [underlined], Ext^1, Ext^1 [underlined] pour des types simples de groupes
% alg. commut. » — each in his hand on a tan / orange card, no \page marker,
% his words the section heading, with a \note. The cover of 159 announces
% the tables of batch 9 (pp. 161-164); p. 160, which sits between, is on
% Ext with supports and does not belong to the tables.
%
% Published text in another hand: pp. 142-144, three leaves paginated 1, 2, 3,
% an English manuscript whose wording coincides, as far as it was read, with
% J. Tate's « Duality theorems in Galois cohomology over number fields »
% (§1 Notation and terminology, §2 Local results, Th. 2.1-2.3, Remark on
% Shatz, §3 Global results, Th. 3.1, Corollaries 1-2, Prop. 3.2, Cor. 1); the
% identification was not checked against the print. One French \note each,
% not re-set; no ink intervention of his was found on them.
%
% Pages summarised: 142-144 as above. None of his own pages is summarised.
% The prose of pp. 147-155 is fast drafting and carries a heavy apparatus;
% p. 157, a sheet of scattered diagrams, is given block by block, its layout
% said in \note{}s; the exact sequence of p. 145 is set line by line.
%
% His pagination: none on his own notes. The Tate leaves carry « 1. », « 2. »,
% « 3. » (top right). No numbered run of his; the brace « 1) », « 2) » of the
% two canonical homomorphisms on p. 160, and the cases (i) of Tate's text.
% Pages 147-155 read as one continuous run (each page begins mid-sentence
% where the last ends); p. 157 is its sheet of figures. No internal page
% reference.
%
% What the batch is about. P. 145 (under the cover « Dualité de Tate-Cassels »,
% after Tate's manuscript): the nine-term Poitou-Tate sequence rewritten for
% O and M, H^r(O, M) ≃ prod H^r(k_p, M) for r > 2, and « If Sha(O, M) = Ker(...)
% is finite ». Pp. 147-155 (under « Groupes de Galois comme groupes analytiques
% sur Q_l, d'après Serre »): « Modules l-adiques sur S » — objects of Ind G(S)
% « de type Q_l/Z_l », of Pro G(S) « de type Z_l », of Ind Pro G(S) « de type
% Q_l »; the categories V_l(S), M_l(S), I_l(S) and their description by
% continuous representations of pi_1(S, xi) when S is connected; base change
% S' -> S; a projective system of preschemes S_i with S = lim S_i and V_l(S) =
% Pseudo lim V_l(S_i) (examples: T_l of an abelian scheme, R^i f_* Z_l); the
% images G_j of pi_j in Aut(V), closed subgroups of finite index in one
% another when the S_i are normal and the transition maps dominant of finite
% type; hence an l-adic Lie algebra g ⊂ gl(V_xi) independent of i, « l'Algèbre
% de Lie de Serre »; its invariance under dominant base extension, contrasted
% with the Lie algebra g' of the image of pi_1; and an Ind-Pro-scheme of Lie
% algebras acting on V. P. 157: figures for the same (pi_1 = lim pi_1(S_i),
% field towers, Gamma ⊂ G x G', and a lemma: H ⊂ G with every element of G
% having a power in H implies H of finite index). P. 160 (under the cover
% « Tableau des Hom ... »): « Cas des préschémas loc. noeth. » — Ext^i_Y(F, G)
% as limits of Ext^i(F/J^{n+1}F, G), a theorem (2) isomorphism for F
% coherent, 1) when X is noetherian), and the start of its proof.
%
% Notation fixed once for the batch: \mathcal{V}_\ell, \mathcal{M}_\ell,
% \mathcal{I}_\ell for his script V, M, I; \mathcal{G}(S) for his script G
% (category of finite étale group schemes over S); \mathcal{S} for his script S
% (the projective system); \mathfrak{g}, \mathfrak{gl} for his gothic g, gl;
% \mathbf{G} for the outlined G of p. 157; \underline{\mathrm{III}} for the
% underlined triple bar (Sha) of p. 145; \coprod for his barred II (direct
% sum), \prod' where he draws the restricted product.
%
% Readings to check first: the exponents of p. 145 (hook-shaped, read 2); the
% two lines of prose at its foot; « Kossul » on p. 152; the first line of
% p. 149 and of p. 152; the margin notes of pp. 150, 153, 154; the struck
% sentence at the top of p. 155; the whole of p. 157's lower third.
% Apparatus: 77 \ill{} and 69 \uncertain{} in the body.

\documentclass[11pt,a4paper]{article}
% Le préambule commun à toutes les transcriptions du fonds Grothendieck.
%
% Il tient en une idée : **l'incertitude se compose, elle ne se résout pas.**
% Une transcription qui lisse un mot douteux a détruit la seule chose pour
% laquelle on la faisait — un lecteur doit pouvoir savoir, à chaque endroit,
% ce qui était sur le feuillet et ce qui a été ajouté. Les six macros qui
% suivent sont donc l'appareil critique, et non de la décoration.
%
% Les mêmes commandes sont reconnues par `scripts/render.mjs`, qui produit la
% vue de lecture du site. Une macro ajoutée ici doit l'être aussi là-bas,
% faute de quoi le rendu échouera bruyamment — ce qui est voulu : un rendu
% silencieusement faux est pire qu'un rendu absent.

\ProvidesPackage{grothendieck}[2026/08/08 Transcriptions du fonds Grothendieck]

\usepackage{amsmath,amssymb,amsthm}
\usepackage{mathrsfs}
\usepackage{stmaryrd}
% Les diagrammes commutatifs sont partout dans ce fonds : c'est la langue
% même de la géométrie algébrique de Grothendieck, et une transcription qui
% les rendrait en prose ne transcrirait rien.
\usepackage{tikz-cd}
% Français par défaut — c'est la langue des feuillets — mais l'anglais est
% chargé aussi : la traduction et le résumé sont des documents anglais, et la
% typographie française (espaces avant les deux-points) y serait une faute.
% Ces documents basculent par \selectlanguage{english} en tête de corps.
\usepackage[english,french]{babel}
\usepackage{fontspec}
\usepackage{geometry}
\geometry{a4paper,margin=25mm,marginparwidth=28mm}
\usepackage{marginnote}
\usepackage{xcolor}
% ulem, not soul. `soul`'s \st and \ul are built on a character-by-character
% scan that XeLaTeX's font machinery breaks: the document fails with an
% undefined control sequence rather than degrading. ulem does the same two jobs
% and is Unicode-engine safe. `normalem` stops it hijacking \emph, which the
% transcriptions use for Grothendieck's own emphasis.
\usepackage[normalem]{ulem}
\usepackage{hyperref}
\hypersetup{colorlinks=true,linkcolor=black,urlcolor=black,pdfborder={0 0 0}}

% --- Métadonnées du lot -----------------------------------------------------
% Reprises telles quelles par le script de rendu, qui les affiche en tête de
% la vue de lecture. Elles fixent ce que le fichier prétend couvrir : sans
% elles, une transcription flotte sans point d'attache dans un fonds de seize
% mille feuillets.

\newcommand{\folder}[1]{\gdef\grothFolder{#1}}
\newcommand{\batch}[1]{\gdef\grothBatch{#1}}
\newcommand{\leaves}[2]{\gdef\grothFirst{#1}\gdef\grothLast{#2}}
\newcommand{\foldertitle}[1]{\gdef\grothFoldertitle{#1}}
\gdef\grothFolder{?}\gdef\grothBatch{?}\gdef\grothFirst{?}\gdef\grothLast{?}\gdef\grothFoldertitle{}

% --- Le résumé --------------------------------------------------------------
%
% Il ouvre la lecture modernisée, sous le titre « Résumé », et il est composé
% en petit. Ce n'est pas un ornement : c'est le seul passage du document qui
% ne soit pas des mathématiques, et un lecteur doit pouvoir le reconnaître
% avant de l'avoir lu — puis le sauter s'il connaît déjà le sujet, ou s'y
% arrêter s'il ne le connaît pas. Le filet qui le suit marque la reprise du
% corps.

\newenvironment{resume}
  {\small\setlength{\parskip}{0.45em}}
  {\par\medskip\hrule\medskip}

% --- Mots-clés ---------------------------------------------------------------
%
% En anglais, en clôture du résumé de la lecture modernisée : le vocabulaire
% moderne sous lequel chercher ce que les feuillets construisent. Le manifeste
% du site (`npm run manifest`) les extrait pour étiqueter la cote entière —
% cette ligne est la source unique des tags, il n'y a pas de fichier à côté.
\newcommand{\keywords}[1]{\par\smallskip\noindent
  {\footnotesize\textsc{Keywords} --- \textit{#1}}\par}

% --- L'appareil critique ----------------------------------------------------

% Le numéro de feuillet, en marge. C'est l'ancre qui permet de régler un
% désaccord sur une formule en regardant *un* feuillet plutôt qu'une chemise
% de six cents. Le site s'en sert aussi pour tourner les pages du fac-similé
% quand on fait défiler la transcription.
\newcommand{\leaf}[1]{%
  \marginnote{\footnotesize\color{gray!70}\textsf{#1}}%
  \label{leaf:#1}\ignorespaces}

% Illisible : on ne devine pas. Un mot inventé qui se lit comme les autres est
% le pire résultat possible.
\newcommand{\ill}{\textcolor{red!60!black}{[\,\ldots\,]}}

% Lecture douteuse : le mot est donné, et signalé comme incertain.
\newcommand{\uncertain}[1]{\uline{#1}}

% Ajout de l'éditeur — une lettre manquante, un mot sous-entendu.
\newcommand{\add}[1]{\textcolor{blue!50!black}{[#1]}}

% Biffé par Grothendieck lui-même. Ce qu'il a rayé fait partie du document :
% une rature dit souvent où la pensée a changé de direction.
\newcommand{\struck}[1]{\sout{#1}}

% Note du transcripteur, sur l'état matériel du feuillet.
\newcommand{\note}[1]{\par\smallskip{\small\color{gray!60!black}\textit{[#1]}}\par\smallskip}

% Note marginale de Grothendieck, distincte du corps du texte.
\newcommand{\marginal}[1]{\par\smallskip\noindent\hspace{1em}%
  {\small\color{gray!60!black}#1}\par\smallskip}

% --- Titre ------------------------------------------------------------------

\AtBeginDocument{%
  \begin{center}
    {\large\bfseries Fonds Alexandre Grothendieck --- cote n\textsuperscript{o}~\grothFolder}\\[2pt]
    {\itshape\grothFoldertitle}\\[4pt]
    {\small feuillets \grothFirst--\grothLast\ (lot \grothBatch)}\\[2pt]
    {\footnotesize\color{gray!60!black}Transcription établie d'après le fac-similé numérisé
     par l'Université de Montpellier.\\
     Les lectures incertaines sont soulignées, les illisibles notées [\,\ldots\,],
     les ajouts entre crochets.}
  \end{center}
  \vspace{1em}}

\endinput


\folder{43}
\batch{8}
\pages{141}{160}
\dating{[à partir de 1961-vers 1968]}
\watermark{Édition de démonstration}
\foldertitle{Dualité des groupes. Jacobiennes généralisées etc : notes et copies de notes manuscrites (s.d.)}

\begin{document}

\section*{Dualité de Tate-Cassels}
\note{titre inscrit de sa main en haut à droite d'une chemise de papier
kraft (page 141), sans autre contenu. Suivent trois feuillets d'une autre
main (pages 142--144), puis sa page 145.}

\page{142}
\note{feuillet manuscrit d'une autre main que la sienne, en anglais, encre
bleu-violet, paginé « 1. » ; texte suivi et soigné, non recomposé ici. Son
libellé coïncide, pour autant qu'il a été lu, avec l'exposé de J.~Tate
« Duality theorems in Galois cohomology over number fields »
(identification non vérifiée sur l'imprimé). Contenu : « §1 Notation and
terminology » --- $A$ anneau de Dedekind, $\bar{A}$ son extension non
ramifiée maximale, $G_{A}$ le groupe de Galois (groupe fondamental de
$\mathrm{Spec}\,A$), $M$ un $G_{A}$-module fini d'ordre $m$ inversible dans
$A$ ; $H^{r}(A, M) = \varinjlim H^{r}(G_{B/A}, M)$ (cohomologie complète
si $G_{A}$ est fini) ; le dual de Cartier $\tilde{M} =
\mathrm{Hom}_{\mathbb{Z}}(M, \bar{A}^{*})$ et le dual de Pontrjagin
$\hat{H}$. « §2 Local Results » --- $k$ localement compact,
$\mathrm{inv}_{k} : H^{2}(k, \bar{k}^{*}) \hookrightarrow \mathbb{Q}/\mathbb{Z}$,
les accouplements $H^{r}(k, M) \times H^{2-r}(k, \tilde{M}) \to
\mathbb{Q}/\mathbb{Z}$ ; Theorem 2.1 ($H^{r}(k, M)$ fini et $H^{2-r}(k,
\tilde{M}) \simeq \widehat{H^{r}(k, M)}$), le cas $k = \mathbb{R}$ ou
$\mathbb{C}$ ; Theorem 2.2 (caractéristique d'Euler-Poincaré $\chi(k, M) =
1/(A : mA) = |m|_{k}$). Aucune intervention de sa main.}

\page{143}
\note{même main, paginé « 2. », non recomposé. Theorem 2.3 (si $mA = A$,
$\mathrm{Card}\,H^{0}(A, M) = \mathrm{Card}\,H^{1}(A, M)$, $H^{2}(A, M) = 0$,
et $H^{1}(A, M)$, $H^{1}(A, \tilde{M})$ sont orthogonaux l'un de l'autre) ;
une « Remark » : le théorème 2.1 généralisé par « S. Shatz (Harvard Thesis,
June 1962) » au cas d'un schéma en groupes artinien commutatif, avec
cohomologie relative aux extensions finies quelconques, groupes munis d'une
topologie localement compacte, et les exemples $\mu_{p}$, $\alpha_{p}$,
$\mathbb{Z}/p\mathbb{Z}$ en caractéristique $p$ (un mot biffé dans le
texte, de la même main) ; « §3 Global results » --- $k$ corps de nombres
(cas (N)) ou de fonctions (cas (F)), $S$ ensemble de places, $A$ l'anneau
des $S$-entiers ; Theorem 3.1 : la suite exacte à neuf termes
$0 \to H^{0}(A, \tilde{M}) \to \prod H^{0}(k_{\mathfrak{p}}, \tilde{M})
\to \widehat{H^{2}(A, M)} \to H^{1}(A, \tilde{M}) \to \cdots$, bouclée
par $\beta_{1}$, $\delta$, $\alpha_{2}$, $\beta_{2}$, et les isomorphismes
$H^{r}(A, M) \simeq \prod H^{r}(k_{\mathfrak{p}}, M)$ pour $r \geqslant 3$ ;
le signe « $(\tilde{M})$ » en marge gauche est de la même main. Aucune
intervention de sa main.}

\page{144}
\note{même main, paginé « 3. », non recomposé. Suite des notations du
théorème 3.1 (produit, somme directe, produit restreint ; les
$\alpha_{i}$, $\beta_{i}$, $\gamma$, $\delta$) ; Corollary 1 ($S$ fini :
$H^{r}(A, M)$ fini, $\chi(A, M)$ et la formule du produit
$\chi(A, M)\chi(A, \tilde{M})/\mathrm{Card}\,H^{3}(A, M)$ égale à $1$ en
cas (F), à $m^{-[k:\mathbb{Q}]}$ en cas (N)) ; Corollary 2 (noyau de
$H^{2}(A, M) \to \coprod H^{2}(k_{\mathfrak{p}}, M)$ fini et dual du noyau
de $\alpha_{1}(\tilde{M})$) ; la construction finie de ce noyau, et
Proposition 3.2 (image par inflation du noyau de
$H^{1}(G, M) \to \prod H^{1}(G_{i}, M)$), Corollary 1 (annulé par le
quotient du degré global par le p.p.c.m. des degrés locaux). Aucune
intervention de sa main.}

\page{145}
\note{au crayon, sur papier kraft ; la suite occupe un grand cadre de trois
lignes. Les exposants tracés en crochet, semblables à un « ? », sont lus
$2$ ; le symbole sous lequel est écrit $p \in S$ est tantôt un II barré
(rendu $\coprod$), tantôt un II gras (rendu $\prod'$).}
\[
0 \to \widehat{\coprod_{p \in S} H^{2}(k_{p}, \tilde{M})} \to \widehat{H^{2}(O, \tilde{M})} \to \prod_{p \in S} \hat{H}^{0}(k_{p}, M) \to \widehat{H^{2}(O, \tilde{M})}
\]
\note{au-dessus de la flèche entre le deuxième et le troisième terme, un
renvoi en V insère « $H^{0}(O, M)$ », surligné ; le chapeau du
$\hat{H}^{0}(k_{p}, M)$ est une petite ligne ondulée. La lecture
« $\coprod$ » du premier terme est incertaine.}
\[
\to H^{1}(O, M) \to \prod_{p \in S}{}' H^{1}(k_{p}, M) \to \widehat{H^{2}(O, \tilde{M})}
\]
\note{l'exposant du dernier terme est de même forme que les précédents et
lu $2$ ; dans la suite de Tate (page 143) le terme à cette place est
$\widehat{H^{1}}$.}
\[
\to H^{2}(O, M) \to \coprod_{p \in S} H^{2}(k_{p}, M) \to 0
\]
\[
\boxed{H^{r}(O, M) \simeq \prod_{p \in S}{}' H^{r}(k_{p}, M)} \qquad r > 2
\]

If
\[
\underline{\mathrm{III}}(O, M) = \mathrm{Ker}\Bigl(H^{1}(O, M) \to \prod_{p \in S} H^{1}(k_{p}, M)\Bigr)
\]
is finite
\note{« If » et « is finite » sont en anglais ; $\underline{\mathrm{III}}$
rend son III souligné (le groupe de Tate-Šafarevič). L'exposant de
$H^{1}$ est peu net.}

$M$ \ill{} dualité sur un \uncertain{corps de nombres} ---

\section*{Groupes de Galois comme groupes analytiques sur $\mathbb{Q}_{\ell}$, d'après Serre}
\note{titre inscrit de sa main en haut à droite d'une chemise de papier
kraft (page 146), sans autre contenu. Les pages 147--155 forment une
rédaction suivie, à l'encre bleue ; la page 157 en porte les figures.}

\page{147}
\emph{Modules $\ell$-adiques sur $S$.}

Soit $S$ un préschéma, $\ell$ un premier, \struck{\ill{}}
\struck{\uncertain{Ind} \ill{}} $\mathcal{G}(S)$ la catégorie des schémas en
groupes (finis) étales sur $S$. Un objet de $\mathrm{Ind}(\mathcal{G}(S))$
est dit « de type $\mathbb{Q}_{\ell}/\mathbb{Z}_{\ell}$ » s'il est
\emph{strict}, et si ses fibres géométriques sont des limites inductives
isomorphes à $(\mathbb{Q}_{\ell}/\mathbb{Z}_{\ell})^{\nu}$ ($\nu$
convenable). Si $S$ est connexe, \struck{$\xi$ un pt géométrique de $S$,
dans} à $\simeq$ près \uncertain{est} indépendant du choix du pt
géométrique de $S$. Si $\xi$ est un tel pt géométrique, la
\struck{\ill{}} catégorie des $G$ « de type
$\mathbb{Q}_{\ell}/\mathbb{Z}_{\ell}$ » est \struck{isomorphe} équivalente
à la catégorie des $\pi_{1}(S, \xi)$-modules discrets qui sont isomorphes
comme groupes abéliens à $(\mathbb{Q}_{\ell}/\mathbb{Z}_{\ell})^{\nu}$,
\struck{\ill{}} ($\nu$ convenable) i.e.\ dont le dual au sens de
Pontrjagin (à valeurs dans $\mathbb{Q}/\mathbb{Z}$) est iso à un module
libre de type fini sur $\mathbb{Z}_{\ell}$.
\note{« (finis) » est ajouté en marge, entouré ; « à $\simeq$ près
\uncertain{est} indépendant du » est écrit au-dessus de la partie biffée de
la ligne, et la phrase n'a pas de sujet exprimé.}

Par dualité (de Pontrjagin) on obtient les objets de
$\mathrm{Pro}\,\mathcal{G}(S)$ qui sont « de type $\mathbb{Z}_{\ell}$ ».
Un objet de $\mathrm{Ind}\,\mathrm{Pro}\,\mathcal{G}(S)$ est dit

\page{148}
« de type $\mathbb{Q}_{\ell}$ » s'il est isomorphe à un \struck{\ill{}}
système inductif $(G_{n})_{n \in \mathbb{N}}$, où pour tout $n$, on a
$G_{n} = G_{0}$, $G$ \uncertain{étant} « de type \struck{$\mathbb{Q}$}
$\mathbb{Z}_{\ell}$ », et que \uncertain{$\varphi_{mn}$} (pour $m
\geqslant n$) est la multiplication par $\ell^{m-n}$. Si $S$ est connexe,
de pt géométrique $\xi$, alors la catégorie des ces $G$ est équivalente à
la catégorie des espaces vectoriels $V$ de dim finie sur
$\mathbb{Q}_{\ell}$, munis d'un homom.\ continu de groupes topologiques
\[
\pi_{1}(S, \xi) \longrightarrow \mathrm{Aut}(V).
\]
\note{un mot biffé devant $\mathrm{Aut}$.}
[Notons que $\pi_{1}(S, \xi)$ étant compact, il y a \ill{} un « réseau »
invariant \ldots]. Soit de façon générale $\mathcal{V}_{\ell}(S)$ la
catégorie des objets de $\mathrm{Ind}\,\mathrm{Pro}\,\mathcal{G}(S)$ qui
sont « de type $\mathbb{Q}_{\ell}$ », $\mathcal{M}_{\ell}(S)$ la catégorie
des \emph{objets} de $\mathrm{Pro}\,\mathcal{G}(S)$ qui sont « de type
$\mathbb{Z}_{\ell}$ », \struck{\ill{} \ill{} changement de \ill{}} \struck{\ill{}} alors
\[
M \rightsquigarrow M \otimes_{\mathbb{Z}_{\ell}} \mathbb{Q}_{\ell} = \bigl(M_{n} = M,\ M_{m} \xrightarrow{\ \ell^{m-n}\ } M_{n}\bigr)
\]
est un foncteur
\[
\mathcal{V}_{\ell}(S) \longrightarrow \mathcal{M}_{\ell}(S)
\]
qui est essentiellement surjectif, \uncertain{et fidèle}, (pas pleinement
fidèle en général). Les
\note{la flèche du foncteur est tracée de $\mathcal{V}_{\ell}(S)$ vers
$\mathcal{M}_{\ell}(S)$, à l'inverse du sens de $M \rightsquigarrow M
\otimes \mathbb{Q}_{\ell}$ ; transcrite telle quelle. « et fidèle » est
ajouté au-dessus de la ligne et souligné. Une ligne biffée après « de type
$\mathbb{Q}_{\ell}$ », et un mot biffé en début de ligne avant
« alors ».}

\page{149}
\ill{} qui \ill{} \ill{} \uncertain{intéressants} \uncertain{sont} les
$\mathcal{V}_{\ell}(S)$, mais en pratique on \uncertain{travaille}
\uncertain{souvent} avec des éléments de $\mathcal{M}_{\ell}(S)$ [ou de façon
équivalente par dualité de Pontrjagin, les $\mathcal{I}_{\ell}(S)$].

Soit $S' \to S$, il s'en déduit un diagramme de foncteurs naturellement
commutatif
\note{« diagramme de » est ajouté au-dessus de la ligne.}

\begin{tikzcd}
\mathcal{I}_{\ell}(S)^{\circ} \arrow[r, "\approx"] \arrow[d] & \mathcal{M}_{\ell}(S) \arrow[r] \arrow[d] & \mathcal{V}_{\ell}(S) \arrow[d] \\
\mathcal{I}_{\ell}(S')^{\circ} \arrow[r, "\approx"] & \mathcal{M}_{\ell}(S') \arrow[r] & \mathcal{V}_{\ell}(S')
\end{tikzcd}

les flèches verticales étant des flèches de changement de base.
\uncertain{Lorsque} $S$ et $S'$ sont connexes, $\xi'$ un pt géométrique de
$S'$, $\xi$ son image sur $S$, d'où
\[
\pi' = \pi_{1}(S', \xi') \longrightarrow \pi = \pi_{1}(S, \xi)
\]
alors les foncteurs verticaux s'interprètent [en termes de
\struck{\ill{}} Modules sur $\pi$ et $\pi'$] par restriction des scalaires.

\page{150}
Supposons que l'on ait un système projectif $\mathcal{S} =
(S_{i})_{i \in I}$, avec $i_{0} \in I$ tel que pour $i \geqslant i_{0}$,
les $\varphi_{i_{0}i} : S_{i} \to S_{i_{0}}$ soient \uncertain{affines}.
Soit $S = \varprojlim S_{i}$. Soit $G$ sur $S$, [d'un des trois types
envisagés]. Dans beaucoup de situations, on désire se ramener au cas d'un
gpe sur $S_{i}$, mais on \uncertain{dispose} \ill{} \ill{} \ill{}
\[
\mathcal{V}_{\ell}(\mathcal{S}) = \text{Pseudo}\varinjlim \mathcal{V}_{\ell}(S_{i})
\]
[resp.\ de $\mathcal{V}_{\ell}$ remplacé par $\mathcal{M}_{\ell}$,
$\mathcal{I}_{\ell}$].

Exemple : Si \uncertain{quasi-compacts}, on part d'un \struck{\ill{}}
$S$-schéma abélien $A$, d'où pour $i$ grand un $S_{i}$-schéma abélien
$A_{i}$, d'où des $T_{\ell}(A_{i}) \in \mathcal{M}_{\ell}(S_{i})$, qui
définissent un élément de \struck{\ill{}} $\mathcal{M}_{\ell}(\mathcal{S})$
[\ill{}]. Supposons pour simplifier les $S_{i}$ connexes, donc $S$
connexe, soit $\xi$ un pt
\marginal{Autre exemple : $X$ schéma lisse \struck{\ill{}} et propre sur
$S$, \uncertain{provenant} \ill{} \ill{} de $S_{i}$, on \struck{\ill{}}
\uncertain{prend} les faisceaux de Weil
$\varprojlim R^{i}f_{*}(\mathbb{Z}/\ell^{n}\mathbb{Z}) =
R^{i}f_{*}(\mathbb{Z}_{\ell})$ \ill{} \uncertain{torsion} \ldots}
\note{la note marginale est écrite en oblique dans la marge gauche, en
regard de l'exemple ; le « $A$ » de « schéma abélien $A$ » est ajouté
au-dessus de la ligne.}

\page{151}
géométrique de $S$, $\xi_{i}$ son image dans $S_{i}$, donc
\[
\pi = \pi_{1}(S, \xi) \xrightarrow{\ \sim\ } \varprojlim_{i} \pi_{i} \qquad \pi_{i} = \pi_{1}(S_{i}, \xi_{i}).
\]
Donc un élément (\uncertain{discret}) de $\mathcal{V}_{\ell}(\mathcal{S})$
\struck{s'identifie \ill{} essentiellement \ill{}} est défini par un
vectoriel $V_{i}$ sur $\mathbb{Q}_{\ell}$ où $\pi_{i}$ opère
« continûment », avec la « proviso » que pour un $W_{j}$ analogue sur
$\pi_{j}$, un « homomorphisme » de $V$ dans $W$ est un élément de
\[
\varinjlim_{\substack{k \\ (k \geqslant i, j)}} \mathrm{Hom}_{S_{k}}\bigl(V_{i} \times_{S_{i}} S_{k},\ W_{j} \times_{S_{j}} S_{k}\bigr)
\]

Supposons maintenant que [\uncertain{pour $i$ grand}] les
\struck{$\pi_{1}(S_{j}, \xi)$} $\pi_{j} \to \pi_{i}$ aient des images
d'indice fini. C'est le cas si les $S_{i}$ sont \emph{normaux}
[\uncertain{intègres}], et les $\varphi_{ij} : S_{j} \to S_{i}$ dominants et
de type fini [définissant des extensions de corps de type fini].
\note{« \uncertain{pour $i$ grand} » est entouré au-dessus de la ligne ;
« \uncertain{intègres} » est écrit au-dessus de « normaux » ; « définissant
des extensions de corps » est ajouté sous la ligne, rattaché à « de type
fini » par une accolade. Un trait vertical en marge gauche longe le
paragraphe.}

[On est ramené au cas des \emph{corps}, où on distingue le cas d'une
extension \uncertain{transcendante}

\page{152}
pure --- ou d'une \ill{} \uncertain{séparable} \uncertain{primaire}, où
$\pi_{j} \to \pi_{i}$ est surjectif --- et d'une extension finie séparable
--- où $\pi_{j} \to \pi_{i}$ est un iso sur un sous-groupe d'indice fini
--- et d'une extension finie radicielle --- où $\pi_{j} \simeq \pi_{i}$
---. Alors les images $G_{j}$ des
\[
\pi_{j} \longrightarrow \mathrm{Aut}\,V_{\xi}
\]
($j \geqslant i_{0}$) sont des sous-groupes fermés, \struck{\ill{}}
d'indice fini les uns dans les autres. Le système projectif des
sous-groupes de $\mathrm{Aut}(V_{i_{0}})$ \struck{\ill{}} $\mathcal{G}$ a
un caractère intrinsèque, est \uncertain{remplacé} par un
\uncertain{équivalent} si on remplace $(S_{i})$ par un système cofinal.
Par \uncertain{Kossul}, \emph{on sait que les sous-groupes analytiques
$\ell$-adiques} de $\mathrm{Aut}(V_{i_{0}})$, dont il est un système, de
Lie
\[
\mathfrak{g} \subset \mathfrak{gl}(V_{\xi}) = \mathrm{End}_{\mathbb{Q}_{\ell}}(V_{\xi})
\]
\emph{indépendante}
\note{sous $\mathrm{Aut}\,V_{\xi}$ un mot biffé ; le nom propre, commençant
par un K, se lit « Kossul » ; il n'a pas été identifié. Le passage « on
sait que les sous-groupes analytiques $\ell$-adiques » est souligné, comme
le dernier mot de la page.}

\page{153}
\emph{de $i$}. \struck{D'\ill{}} + C'est \emph{l'Algèbre de Lie de Serre}
\note{la fin de ligne est soulignée deux fois ; le « Lie » n'est pas sûr,
le mot pouvant se lire « l'Algèbre de Serre ».}

Supposons $S = \mathrm{Spec}(A)$, $A$ \uncertain{normal}
\struck{\uncertain{intègre}} \struck{\ill{}} \struck{\ill{} \ill{} \ill{}}.
Alors on a
\[
A = \varinjlim_{i} A_{i}
\]
les $A_{i}$ \uncertain{sous-anneaux} de type fini ;
\struck{\ill{} normaux sur $\mathbb{Z}$} normaux [\uncertain{lorsque $A$
l'est}] \uncertain{si} on y tient.
[\textbf{N.B.} Si] Si $A$ contient un corps premier, si ce dernier est de
car $p > 0$, les $A_{i}$ sont des algèbres de type fini sur
$\mathbb{F}_{p}$. Si $A \supset \mathbb{Q}$, on peut aussi \ill{}
\ill{}, et écrire
\[
A = \varinjlim_{\alpha} B_{\alpha}
\]
les $B_{\alpha}$ les \uncertain{sous}-$\mathbb{Q}$-algèbres de $A$ qui sont
de type fini ; \struck{Mais \uncertain{plus} \ill{}}
$\mathcal{V}((A_{i}))$ est \uncertain{alors} \uncertain{défini} \ill{}
\ill{} $\mathcal{V}((B_{\alpha}))$, il \uncertain{semble} d'ailleurs que les
objets \struck{\ill{}} de $\mathcal{V}((B_{\alpha}))$ qui \ill{} \ill{}
\uncertain{fabriquer} proviennent en fait de $\mathcal{V}((A_{i}))$, (cf
exemples plus haut). On dira
\marginal{N.B. Pour la \uncertain{question} \ill{}, on \ill{} \emph{des
algèbres de} \ill{} des $A_{i}$ \struck{\ill{}} \ill{} $B_{\alpha}$ \ill{}
\ill{}}
\note{« N.B. Si » est encadré en début de ligne, avec un trait courbe ; la
ligne biffée « \ill{} normaux sur $\mathbb{Z}$ » commence par une lettre
repassée. « lorsque $A$ l'est » est ajouté au-dessus de la ligne,
rattaché par une accolade ; « $\mathbb{Q}/$ » est ajouté au-dessus de
« sous-algèbres » ; au-dessus de « Mais \ldots » biffé, deux mots
illisibles. La note marginale est écrite en oblique dans la marge gauche,
en bas.}

\page{154}
qu'on a un \struck{Module} « vectoriel $\ell$-adique » [ou « Module
$\ell$-adique »] \emph{attaché à $A$} si on a un élément de
\[
\mathcal{V}((A_{i})) = \text{Pseudo}\varinjlim \mathcal{V}(A_{i}).
\]
Si $A$ est normal (\uncertain{p.ex.}), il lui est attaché \struck{\ill{}}
canonique [pour tout pt géométrique $\xi$ de $S$] une algèbre de Lie
$\mathfrak{g}_{\xi}$ de $\mathfrak{gl}(V_{\xi})$. D'ailleurs, l'image de
$\pi_{1}(S, \xi) \to \mathrm{Aut}(V_{\xi})$ est un sous-groupe analytique,
dont une algèbre de Lie $\mathfrak{g}'_{\xi} \subset \mathfrak{g}_{\xi}$. La
formation de $\mathfrak{g}_{\xi}$ est \struck{compatible}
\emph{\uncertain{invariante} par extension (dominante) de la base},
contrairement à celle de $\mathfrak{g}'_{\xi}$ [$\mathfrak{g}'_{\xi}$
\emph{diminue} quand $(S, \xi)$ est remplacé par $(S', \xi')$
\uncertain{convenable}]. Si d'ailleurs $S$ est de type fini sur \ill{}, on a
évidemment $\mathfrak{g}'_{\xi} = \mathfrak{g}_{\xi}$. On constate
facilement que [dans le cas général d'un système $(S_{i})$ \ldots]
\struck{\ill{}} le choix de « \uncertain{classes} de \uncertain{chemins} »
de $\xi$ à $\xi'$ \struck{\ill{}} (pts géom.\ de $S$) \uncertain{définit}
\marginal{Plus \uncertain{commode} : $\mathfrak{g}_{\xi}$ qui \ill{}
\uncertain{extension} de la base par \ill{} dominantes}
\note{« (p.ex.) » est écrit au-dessus de « normal », rattaché par une
accolade ; « canonique » remplace au-dessus de la ligne un mot biffé ;
« (dominante) » est ajouté au-dessus de « extension ». La note marginale,
encadrée, est reliée par un signe ¶ à la fin du crochet
« [$\mathfrak{g}'_{\xi}$ diminue \ldots] ».}

\page{155}
un iso $V_{\xi} \to V_{\xi'}$ qui transporte $\mathfrak{g}_{\xi}$ en
$\mathfrak{g}_{\xi'}$. \struck{En particulier, $\mathfrak{g}_{\xi}$ est
invariant sous $\pi_{1}(S_{i}, \xi)$} ; [\uncertain{tous} \ill{} \ill{}
\ill{} \ill{} (\uncertain{formules} isomorphismes)] --- on a un foncteur qui
à tout « vectoriel de type $\mathbb{Q}_{\ell}$ attaché à $(S_{i})$ »,
[\uncertain{sur} $V$,] associe une « algèbre de Lie de type
$\mathbb{Q}_{\ell}$ attachée à $(S_{i})$ », \struck{dite} [\uncertain{sur}
$V$] $\mathfrak{g}$, \struck{qui \ill{} « algèbre de} [disons une
Ind-Pro-schéma en algèbres de Lie sur $S$ convenable] \uncertain{opérant}
sur $V$ de \uncertain{façon} évidente.
\note{le trait qui barre « En particulier \ldots\ $\pi_{1}(S_{i}, \xi)$ »
peut n'être qu'un soulignement ; la parenthèse en dessous est rattachée
par une accolade à « qui à tout ». « sur $V$ » est ajouté deux fois au-dessus
de la ligne.}

\page{157}
\note{feuillet de figures éparses, à l'encre bleue, relatif aux pages
147--155 ; donné bloc par bloc, de haut en bas.}
\[
\pi_{1}(S, \xi) = \varprojlim \pi_{1}(S_{i}, \xi_{i}) \longrightarrow \mathbf{G} \qquad \mathfrak{g}
\]
\note{la flèche descend en oblique vers $\mathbf{G}$ (son G ajouré) ; le
$\mathfrak{g}$ est isolé en haut à droite, avec un arc de cercle au-dessous.
À gauche de $\mathbf{G}$, isolé : « $\pi_{1}(S, \xi)$ ».}

\begin{tikzcd}
 & \bar{K}_{2} \arrow[d, no head] \\
\bar{K}_{1} \arrow[r, no head] \arrow[d, no head] & (\bar{K}_{1}, K_{2}) \arrow[d, no head] \\
K_{1} \arrow[r, no head] \arrow[d, no head] & K_{2} \\
k_{1} &
\end{tikzcd}

\note{sous le trait $K_{1}$ --- $K_{2}$, un mot entre guillemets,
illisible ; le $k_{1}$ du bas est \uncertain{minuscule}.}

$\exists\, g \in G$, $U$ \uncertain{voisinage} de $e$, $(gu)^{n} \in H$
\uncertain{pour tout} $u$
\note{à gauche du treillis.}
\[
A \to A_{i}, \quad A' \qquad S \qquad A \to B, \quad S \leftarrow T, \quad S_{i} \leftarrow T_{i}, \quad \pi_{1}(S_{i}) \leftarrow \pi_{1}(T)
\]
\note{$A'$ est au-dessus de $A_{i}$, relié à $A$ par un trait ; à droite,
un $S$ au-dessus d'un trait vertical double. Les trois flèches $S \leftarrow
T$, $S_{i} \leftarrow T_{i}$, $\pi_{1}(S_{i}) \leftarrow \pi_{1}(T)$ sont
empilées sous $A \to B$.}
\[
S \leftarrow S' \qquad \mathbf{G} \longleftrightarrow \mathbf{G}' \qquad \Gamma \subset G \times G' \qquad G \to G'
\]
\note{de $\Gamma$ partent deux flèches, vers $G$ et vers $G'$ ; un trait
oblique épais précède $\mathbf{G} \leftrightarrow \mathbf{G}'$. À droite,
empilés : « $\mathfrak{h}_{1} \subset \mathfrak{h}_{2}$ », « $H_{1}
\subset$ \struck{\ill{}} », « $H_{1} \subset H_{2}$ », avec $\cap$ et $G$
sous $H_{2}$.}

\uncertain{$\ell$-gpes} de type \ill{} \qquad \ill{} groupes
\uncertain{finis}
\[
\Updownarrow
\]
systèmes projectifs d'algèbres de Lie de dim finie sur $\mathbb{Q}_{\ell}$
\note{la double flèche est tracée comme une flèche montante biffée à côté
d'une flèche descendante ; à droite, dans une ellipse, « \uncertain{Fidèle}
\ill{} ».}

\emph{Lemme} Le th.\ est vrai pour $H_{1} \simeq \mathbb{Z}_{\ell}$

que $H_{1} \subset$ \struck{$G$} $H_{2} \Rightarrow \mathfrak{h}_{1} \subset
\mathfrak{h}_{2} \overset{?}{\Longrightarrow}$ $H$ \struck{\ill{}}
\uncertain{engendré} du ss-groupe où une puissance de $H_{1}$ \ill{}
\struck{\ill{}} \struck{$G$ \ill{} $H_{1}$ \ill{}} \ill{} dans $H_{2}$
$\Rightarrow$ tout élément de $H_{1}$ a une puissance dans $H_{2}$
\note{sous $\mathfrak{h}_{1} \subset \mathfrak{h}_{2}$, à gauche :
« i.e.\ $H_{1} \cap H_{2}$ est ouvert dans $H_{1}$, i.e.\ d'indice fini » ;
un arc relie le « ? » au Lemme au-dessus. Un trait courbe au crayon sépare
ce bloc des systèmes projectifs d'algèbres de Lie.}

par Lemme : Si $H \subset G$, et tout élément de $G$ a une puissance dans
$H$, alors $H$ d'indice fini dans $G$.
\note{écrit au bas de la page, avec une longue flèche qui revient vers
« i.e.\ d'indice fini ».}

\section*{Tableau des Hom, $\underline{\mathrm{Hom}}$, $\mathrm{Ext}^{1}$, $\underline{\mathrm{Ext}}^{1}$ pour des types simples de groupes alg.\ commut.}
\note{titre inscrit de sa main en haut à droite d'une chemise de papier
orangé (page 159), sans autre contenu ; « Hom » et « Ext$^{1}$ » soulignés
la seconde fois. Il annonce les tableaux des pages 161--164 (lot 9) ; la
page 160, qui s'intercale, porte sur les Ext à supports.}

\page{160}
\emph{Cas des préschémas} loc.\ noeth.

Si $F$ est cohérent, les $\underline{\mathrm{Ext}}^{i}_{Z}(F, G)$ sont
quasi-cohérents (cela se voit par exemple sur la suite spectrale de terme
initial $\underline{H}^{p}_{Z}(\underline{\mathrm{Ext}}^{q}_{\mathcal{O}_{X}}(F, G))$,
utilisant le fait que \uncertain{tout} extension de quasi-cohérents est
\uncertain{idem} \ldots

On a des homom.\ canoniques (\struck{$Z$} fermé)
\begin{align*}
&1) \quad \mathrm{Ext}^{i}_{Y}(X; F, G) \longleftarrow \varinjlim \mathrm{Ext}^{i}(X; F/\mathcal{J}^{n+1}F, G) \\
&2) \quad \underline{\mathrm{Ext}}^{i}_{Y}(F, G) \longleftarrow \varinjlim \underline{\mathrm{Ext}}^{i}(F/\mathcal{J}^{n+1}F, G)
\end{align*}
\note{les deux lignes sont réunies par une accolade à gauche ; dans 1),
l'indice $Y$ est écrit sur un $Z$ biffé, et dans la parenthèse le $Z$ est
biffé sans lettre qui le remplace. L'idéal est son J cursif.}

\emph{Th.} [Supposons $F$ \emph{cohérent}.] L'homom 2) est un isom,
\struck{si $X$ \ill{}}, et 1) est un isom si $X$ est noethérien.
\note{« Supposons $F$ cohérent. » est ajouté au-dessus de la ligne,
rattaché par une accolade.}

On déduit \struck{\ill{}} la \uncertain{deuxième} assertion de la
\uncertain{première} comme d'habitude, par la suite spectrale
\[
H^{p}(X, \underline{\mathrm{Ext}}^{q}_{Y}(F, G)) \Longrightarrow \mathrm{Ext}^{*}_{Y}(X; F, G).
\]
L'énoncé est de nature \emph{locale}, on peut supposer $X$ affine. Pour
$i = 0$, c'est évident, il suffit de \struck{examiner} prouver que si $G$
\uncertain{provient} d'un \uncertain{module} \emph{injectif}, on a
$\underline{\mathrm{Ext}}^{i}_{Y}(F, G) = 0$ pour $i > 0$. Or on regarde
$\mathrm{Ext}^{p}_{A}(M, H^{q}_{J}(N))$ et on note que si $N$
\note{la phrase s'interrompt au bas de la page ; la page 161 (lot 9)
commence le tableau des Hom.}

\end{document}
